\documentclass[11pt]{article}
\usepackage[a4paper,margin=26mm]{geometry}
\usepackage{amsmath,amssymb,amsthm,booktabs,tabularx,array}
\usepackage[T1]{fontenc}
\usepackage{lmodern}
\usepackage[colorlinks=true,allcolors=blue,breaklinks=true]{hyperref}
\newtheorem{theorem}{Theorem}
\newtheorem{corollary}[theorem]{Corollary}
\newcommand{\R}{\mathbb R}
\newcommand{\In}{\operatorname{In}}
\newcommand{\Span}{\operatorname{span}}
\newcommand{\diag}{\operatorname{diag}}
\newcommand{\rank}{\operatorname{rank}}
\title{MathOverflow 513954, revision 2:\\
singular signed downdates, fixed coordinates,\\
and a sharp polynomial sphere bound}
\author{Haobo Ma\\\small ChronoAI Pte Ltd, Singapore
\and Wenlin Zhang\\\small National University of Singapore\\
\small ChronoAI Pte Ltd, Singapore}
\date{November 2, 2026}
\begin{document}
\maketitle
\begin{abstract}
MathOverflow 513954, revision 2~\cite{lucid} has affirmative answers:
we give a self-contained proof of the inertia formula previously proved
by user1551~\cite{user1551}, and prove the continuous-family assertion.
For any real symmetric matrix $A$, singular or otherwise, and
$\Delta=x^{\mathsf T}Ax\ne0$, the matrix
$M=A-(Ax)(Ax)^{\mathsf T}/\Delta$ has its positive index decreased by one
when $\Delta>0$, or its negative index decreased by one when $\Delta<0$;
the other nonzero index is unchanged, and the nullity increases by one.
Its kernel is $\ker A\oplus\Span\{x\}$.
For continuous diagonal families with $M\succeq0$ on a nonempty connected
space, $\Delta$ has constant sign; if negative, a unique diagonal coordinate
is negative everywhere and all other coordinates are nonnegative.
The conclusions extend to countably punctured simplices, and a block
congruence removes any nondegenerate subspace. For the linked sphere
question~\cite{otherquestion}, the increasing odd polynomial
$h(t)=(3t-t^3)/2$ has squared ratio infimum
$16(17-9\sqrt3)/23$ over all dimensions, strictly below the
two-coordinate infimum $80/81$.
\end{abstract}
\noindent\small\textcopyright\ 2026 Haobo Ma and Wenlin Zhang.
Licensed under \href{https://creativecommons.org/licenses/by/4.0/}{CC BY 4.0}.
\normalsize

\section{The inertia question}
Lucid's question~\cite{lucid}, revision 2 (August 5, 2026), considers
$D=\diag(d_1,\ldots,d_n)$, $w=Dx$ and $M=D-ww^{\mathsf T}/\Delta$,
where $\Delta=x^{\mathsf T}Dx\ne0$.
We use \emph{signed downdate} for this two-sign formula: with
$u=Dx/\sqrt{|\Delta|}$ it is $D-uu^{\mathsf T}$ when $\Delta>0$ and
$D+uu^{\mathsf T}$ when $\Delta<0$, a rank-one update in the latter case.
Both the inertia formula and the related continuous-family assertion have
affirmative answers. On Lucid's Mathematics Stack Exchange cross-post,
user1551 gave a proof of the matrix inertia formula on August 7,
2026~\cite{user1551}; its zero count follows by subtracting the positive
and negative counts from the dimension. We credit that earlier proof and
give a self-contained completion of Lucid's proposed hyperplane argument,
together with the continuous-family proof.
The proposed decomposition in the question already
contains the matrix argument: an invertible change of basis is required,
but invertibility of $D$ is not. We give the argument for arbitrary real
symmetric $A$ and then prove the fixed-index assertion separately.

For a real symmetric matrix $B$, write
$\In(B)=(n_+(B),n_-(B),n_0(B))$ for the positive, negative and zero
eigenvalue counts, including multiplicity.
Congruence preserves these counts even for a degenerate form.
Indeed, by the spectral theorem, $n_+(B)$ is the largest dimension of a
subspace on which $v^{\mathsf T}Bv$ is positive definite: the positive
eigenspace attains it, and any larger subspace intersects the nonpositive
eigenspace. An invertible substitution preserves this largest dimension.
Apply the same argument to $-B$ for $n_-$, and use
$n_0=n-n_+-n_-$ for the remaining count. This is Sylvester's law of inertia
in precisely the form needed below.

\begin{theorem}\label{thm:matrix}
Let $A\in\R^{n\times n}$ be symmetric and $x\in\R^n$ satisfy
$\Delta=x^{\mathsf T}Ax\ne0$. Put $w=Ax$ and
$M=A-ww^{\mathsf T}/\Delta$. Then
\begin{equation}\label{eq:inertia}
\In(M)=
\begin{cases}
(n_+(A)-1,\ n_-(A),\ n_0(A)+1),&\Delta>0,\\
(n_+(A),\ n_-(A)-1,\ n_0(A)+1),&\Delta<0.
\end{cases}
\end{equation}
Moreover,
\begin{equation}\label{eq:kernel}
\ker M=\ker A\oplus\Span\{x\},\qquad \rank M=\rank A-1.
\end{equation}
\end{theorem}
\begin{proof}
The nonzero pivot implies $w\ne0$ and $x\ne0$. Let
$H=\{z\in\R^n:w^{\mathsf T}z=0\}$. Every $v$ has the unique decomposition
\[
v=\left(v-\frac{w^{\mathsf T}v}{\Delta}x\right)
  +\frac{w^{\mathsf T}v}{\Delta}x\in H\oplus\Span\{x\}.
\]
Choose a basis matrix $Z$ for $H$ and put $P=[Z\ \ x]$.
Thus $P$ is invertible, including when $A$ is singular.
Symmetry gives $Z^{\mathsf T}Ax=Z^{\mathsf T}w=0$, and $Mx=0$.
With $C=Z^{\mathsf T}AZ$ we therefore have
\begin{equation}\label{eq:blocks}
P^{\mathsf T}AP=\begin{pmatrix}C&0\\0&\Delta\end{pmatrix},
\qquad
P^{\mathsf T}MP=\begin{pmatrix}C&0\\0&0\end{pmatrix}.
\end{equation}
Congruence and direct-sum additivity give
$\In(A)=\In(C)+\In([\Delta])$ and
$\In(M)=\In(C)+(0,0,1)$, proving~\eqref{eq:inertia}.
In particular the subtracted positive or negative count is at least one.

For the kernel statement, if $Av=0$ then
$w^{\mathsf T}v=x^{\mathsf T}Av=0$, so $Mv=0$.
Conversely, $Mv=0$ implies
$A(v-tx)=0$ with $t=w^{\mathsf T}v/\Delta$.
Together with $Mx=0$ this gives the sum in~\eqref{eq:kernel}.
If $tx\in\ker A$, then $t\Delta=x^{\mathsf T}A(tx)=0$, hence $t=0$;
the sum is direct. Rank-nullity gives the rank identity.
\end{proof}

For example, $D=\diag(-2,0,3)$ and $x=(1,0,0)^{\mathsf T}$ give
$\Delta=-2$ and $M=\diag(0,0,3)$. The old null direction remains, and the
new one is independent of it. No inverse of $D$ appears anywhere.

\begin{corollary}\label{cor:psd}
Under the hypotheses of Theorem~\ref{thm:matrix},
\[
M\succeq0\quad\Longleftrightarrow\quad
\begin{cases}
A\succeq0,&\Delta>0,\\
n_-(A)=1,&\Delta<0.
\end{cases}
\]
For diagonal $A=D$ in the negative-pivot case, this means exactly one $d_j<0$
and $d_i\ge0$ for all $i\ne j$.
\end{corollary}
\begin{proof}
A symmetric matrix is positive semidefinite exactly when its negative
index is zero. Substitute the corresponding negative count
from~\eqref{eq:inertia}. The eigenvalues of a diagonal matrix are its entries.
\end{proof}

\begin{corollary}[Nondegenerate block removal]\label{cor:block}
Let $A$ be real symmetric, $X\in\R^{n\times k}$, and
$G=X^{\mathsf T}AX$ invertible. Define
$N=A-AXG^{-1}X^{\mathsf T}A$. Then
\[
\In(N)=\bigl(n_+(A)-n_+(G),\ n_-(A)-n_-(G),\ n_0(A)+k\bigr),
\]
\[
\ker N=\ker A\oplus\operatorname{im}X,\qquad
\rank N=\rank A-k.
\]
\end{corollary}
\begin{proof}
The matrix $X$ is injective since $G$ is invertible. With
$H=\ker(X^{\mathsf T}A)$, the identity
$v=(v-XG^{-1}X^{\mathsf T}Av)+XG^{-1}X^{\mathsf T}Av$
gives $\R^n=H\oplus\operatorname{im}X$.
In a basis $P=[Z\ X]$ adapted to this sum,
$P^{\mathsf T}AP=\diag(Z^{\mathsf T}AZ,G)$ and
$P^{\mathsf T}NP=\diag(Z^{\mathsf T}AZ,0_k)$.
Congruence gives the inertia. If $Nv=0$, then
$v-XG^{-1}X^{\mathsf T}Av\in\ker A$; conversely $N$ kills both
$\ker A$ and $\operatorname{im}X$. If $Xc\in\ker A$ then $Gc=0$,
so their intersection is zero. Rank-nullity proves the rank formula.
\end{proof}

\begin{corollary}[Coefficient changes]\label{cor:coefficient}
Under Theorem~\ref{thm:matrix}, put $M_t=A-tww^{\mathsf T}/\Delta$.
Then, componentwise,
\[
\In(M_t)=\In(A)-\In([\Delta])+\In([(1-t)\Delta]).
\]
Thus $t<1$ preserves inertia, $t=1$ gives~\eqref{eq:inertia},
and $t>1$ replaces one eigenvalue sign by its opposite without changing nullity.
\end{corollary}
\begin{proof}
The same basis in~\eqref{eq:blocks} gives
$P^{\mathsf T}M_tP=\diag(C,(1-t)\Delta)$.
\end{proof}

\section{The continuous-family corollary}
\begin{theorem}\label{thm:family}
Let $Y$ be a nonempty connected topological space, and let
$d,x:Y\to\R^n$ be continuous. For $y\in Y$ define
\begin{align*}
D(y)&=\diag(d_1(y),\ldots,d_n(y)),\qquad
\Delta(y)=x(y)^{\mathsf T}D(y)x(y),\\
M(y)&=D(y)-\frac{D(y)x(y)x(y)^{\mathsf T}D(y)}{\Delta(y)}.
\end{align*}
Suppose $\Delta(y)\ne0$ and $M(y)\succeq0$ for every $y$.
Then exactly one of the following alternatives holds:
\begin{enumerate}
\item $\Delta(y)>0$ and every $d_i(y)\ge0$ for all $y$;
\item $\Delta(y)<0$ for all $y$, and there is a unique index $j$ such that
      $d_j(y)<0$ for every $y$ and $d_i(y)\ge0$ for every $i\ne j$ and every $y$.
\end{enumerate}
\end{theorem}
\begin{proof}
The continuous nonzero function
$\Delta(y)=\sum_i d_i(y)x_i(y)^2$ has constant sign:
the inverse images of $(0,\infty)$ and $(-\infty,0)$ are disjoint open
sets covering $Y$, so connectedness and nonemptiness force exactly one
to be empty. If the sign is positive, Corollary~\ref{cor:psd} proves (1).
If negative, that corollary shows that at each point exactly one diagonal
entry is negative. The sets $U_i=\{y:d_i(y)<0\}$ are open, pairwise
disjoint, and cover $Y$. Each is also closed because its complement is
the union of the other $U_i$. A nonempty $U_j$ therefore equals $Y$.
All other $U_i$ are empty, and the asserted uniqueness follows.
\end{proof}

\begin{corollary}\label{cor:simplex}
Let $n\ge3$, let $X$ be homeomorphic to an open $(n-1)$-simplex, and let
$F\subset X$ be finite. With $Y=X\setminus F$, all conclusions of
Theorem~\ref{thm:family} hold under its stated continuity, nonzero-pivot
and positive-semidefiniteness hypotheses.
\end{corollary}
\begin{proof}
Transport $F$ to the finite set $B$ in the convex open simplex
$S=\{u\in\R^{n-1}:u_i>0,\ \sum_i u_i<1\}$.
For $p,q\in S\setminus B$, choose $z\in S$ outside the finitely many
lines through $p,b$ or $q,b$ with $b\in B$.
Such a $z$ exists: a small circle in a two-dimensional affine plane
inside $S$ meets each of these lines in at most two points.
The segments $[p,z]$ and $[z,q]$ stay in $S$ by convexity and avoid $B$
by the choice of $z$. Thus $S\setminus B$ is path connected and is nonempty
because $S$ is open and $B$ finite. The homeomorphism transfers these
properties to $Y$, and Theorem~\ref{thm:family} applies.
\end{proof}

\begin{corollary}[Countably many punctures]\label{cor:countable}
Corollary~\ref{cor:simplex} remains true for countable $F$.
\end{corollary}
\begin{proof}
In the preceding argument, the forbidden lines meet a small circle
in at most countably many points, whereas the circle is uncountable.
The same choice of $z$ therefore works. The open simplex is uncountable,
so its complement of a countable set is nonempty.
\end{proof}

\begin{corollary}[Sharp dimension threshold]\label{cor:sharp}
The dimension threshold $n\ge3$ in Corollary~\ref{cor:simplex} is sharp
for a conclusion uniform over all finite deleted sets.
\end{corollary}
\begin{proof}
For $n=2$, let $X=(-1,1)$ and $F=\{0\}$.
On $y<0$ take $D=\diag(-1,1)$ and $x=(1,0)^{\mathsf T}$;
on $y>0$ take $D=\diag(1,-1)$ and $x=(0,1)^{\mathsf T}$.
These functions are continuous on $X\setminus F$, with $\Delta=-1$ and
$M\succeq0$ everywhere, but the index of the negative diagonal entry switches
between the two components.
\end{proof}
In Theorem~\ref{thm:family} the other, nonnegative entries may vanish.

\section{The linked sphere question}
Lucid's Mathematics Stack Exchange question 5146381~\cite{otherquestion}
asks whether, for continuous $h:(0,1]\to\R$ with $h(t)\ge h(1)t>0$,
the infimum of
\[
R_h(x)=\frac{\sum_i x_i h(x_i)}{\sqrt{\sum_i h(x_i)^2}},
\qquad \sum_i x_i^2=1,
\]
can always be reduced to at most two active coordinates.
We work with strictly positive coordinates and allow their number to vary.
This is the active-coordinate convention: zero coordinates are omitted.
The source specifies no $h(0)$; retaining zeros instead requires specifying
that value. Our polynomial below has $h(0)=0$, so both conventions agree.
\begin{corollary}[Numerator bound and stationarity]\label{cor:stationarity}
Under the source's hypotheses, $N:=\sum_i x_i h(x_i)\ge h(1)$.
If $h$ is differentiable near the active coordinates of a positive unit
vector with each $x_i<1$, its stationary points satisfy
\[
D^2\bigl[h(x_i)+x_i h'(x_i)\bigr]-N h(x_i)h'(x_i)
=2\lambda x_i D^3,\qquad D:=\sqrt{\sum_j h(x_j)^2}.
\]
Continuity alone is insufficient to justify this equation's derivatives.
\end{corollary}
\begin{proof}
The source's termwise proof~\cite{otherquestion} gives
$N\ge h(1)\sum_i x_i^2=h(1)$.
Differentiating $N/D$ gives the left side of the displayed equation
divided by $D^3$; the sphere constraint has gradient $2x\ne0$.
The Lagrange multiplier condition proves the assertion.
\end{proof}

Two earlier observations illuminate the general question:
daw suggested $h(t)=1$ on August 8, 2026~\cite{daw}, and
Exodd used $h(t)=1/t$ in an answer on August 11~\cite{exodd}.
\begin{corollary}[Infima for the earlier examples]\label{cor:examples}
Under the active-coordinate convention, $h=1$ has fixed-dimension
infimum $1/\sqrt d$, and all-dimension infimum zero, whereas
dimensions at most two give infimum $1/\sqrt2$.
For $h(t)=1/t$, the infimum is already zero in dimension two.
\end{corollary}
\begin{proof}
For $h=1$, $R_h(x)=(\sum_i x_i)/\sqrt d\ge1/\sqrt d$,
and positive vectors approaching $(1,0,\ldots,0)$ give equality
in the limit. Then let $d\to\infty$.
For $h(t)=1/t$ in dimension two, put $p=x_1^2$;
then $R_h(x)^2=4p(1-p)\to0$ as $p\to0$.
\end{proof}

\begin{theorem}[An increasing odd polynomial]\label{thm:sphere}
Let $h(t)=(3t-t^3)/2$ on $[0,1]$. It is strictly increasing,
$h(1)=1$, and $h(t)/t=(3-t^2)/2$ is strictly decreasing for $0<t\le1$.
Over nonnegative unit vectors of arbitrary finite dimension,
\begin{equation}\label{eq:sphereinf}
\inf R_h(x)^2=\frac{16(17-9\sqrt3)}{23}.
\end{equation}
The infimum is not attained in a finite dimension. Restricting to at most
two nonzero coordinates gives the attained minimum $80/81$ instead.
For example, a vector with squared coordinates $(4/5,1/10,1/10)$ gives
\[
R_h(x)^2=\frac{13689}{13885}<\frac{80}{81}.
\]
The infimum in~\eqref{eq:sphereinf} is approached by one squared coordinate
$a_*=6-3\sqrt3$ and $m$ squared coordinates $(1-a_*)/m$ as $m\to\infty$.
\end{theorem}
\begin{proof}
Strict increase follows from $h'(t)=3(1-t^2)/2>0$ for $0\le t<1$.
Put $p_i=x_i^2$, $s_2=\sum_i p_i^2$, and $s_3=\sum_i p_i^3$.
Then $\sum_i p_i=1$ and a direct expansion gives
\begin{equation}\label{eq:moments}
R_h(x)^2=\frac{(3-s_2)^2}{9-6s_2+s_3}.
\end{equation}
If there are at most two positive $p_i$, write them as $p,1-p$,
allowing endpoints, and put $u=p(1-p)\in[0,1/4]$. Then
\[
R_h(x)^2=\frac{(1+u)^2}{1+9u/4},\qquad
R_h(x)^2-\frac{80}{81}
=\frac{(9u-1)^2}{81(1+9u/4)}.
\]
The value $u=1/9$ is feasible, proving the two-coordinate assertion.
For the three-coordinate vector,
\[
(p_1,p_2,p_3)=(4/5,1/10,1/10),\qquad s_2=33/50,\quad s_3=257/500.
\]
Equation~\eqref{eq:moments} gives the displayed counterexample.
The strict comparison is
\[
\frac{80}{81}-\frac{13689}{13885}=\frac{1991}{81\cdot13885}>0.
\]

For arbitrary dimension, put $a=\sqrt{s_2}\in(0,1]$.
Since each $p_i\le a$, we have $s_3\le a s_2=a^3$. Thus
\[
R_h(x)^2\ge f(a):=\frac{(3-a^2)^2}{9-6a^2+a^3}.
\]
The denominator is positive on $[0,1]$, and
\[
f'(a)=-\frac{a^2(3-a^2)(a^2-12a+9)}{(9-6a^2+a^3)^2}.
\]
It is negative before $a_*=6-3\sqrt3\in(0,1)$ and positive after it.
Substitution gives $f(a_*)=16(17-9\sqrt3)/23$.
For the stated sequence,
$s_2=a_*^2+(1-a_*)^2/m\to a_*^2$ and
$s_3=a_*^3+(1-a_*)^3/m^2\to a_*^3$, so the lower bound is sharp.
Equality $s_3=a^3$ in a finite vector forces every positive $p_i=a$.
Then $s_2=a\sum_i p_i=a$ and $s_2=a^2$, so $a=1$ and $R_h(x)^2=1$.
Since $f(a_*)<1$, the infimum is not attained.
\end{proof}

This disproves a universal two-coordinate reduction even for an increasing
odd polynomial with decreasing $h(t)/t$. It also gives an exact
one-large-coordinate plus diffuse-tail limit of the kind asked about in
the source. The broader classification for arbitrary $h$, and the
source's numerical observation for an unspecified odd polynomial,
are separate questions.

\section{Source coverage and verification}\label{sec:coverage}
All mathematical assertions above have written proofs.
The two pinned Lean sources~\cite{inertialean,familylean} establish the
following corresponding statements at the commit listed below.
\begin{flushleft}\small
\texttt{0f1f822c673c1903dc6b0a209a23dacebca2aa6c}\\
Common namespace: \texttt{D5.S3.QuadraticForms}\\
\texttt{I}: \texttt{StructuredRankOneInertia}\\
\texttt{F}: \texttt{StructuredRankOneFamily}
\end{flushleft}
\begin{center}\small
\begin{tabularx}{\linewidth}{@{}>{\raggedright\arraybackslash}p{.32\linewidth}>{\raggedright\arraybackslash}X@{}}
\toprule
Statement & Lean declaration or written consequence\\
\midrule
Both inertia triples, any real symmetric $A$ (hence singular $D$)
 & \texttt{I.inertia\_downdate}\\
Kernel and nullity
 & \texttt{I.kernel\_split}, \texttt{I.nullity\_downdate}\\
Punctured-simplex path connectedness
 & \texttt{F.punctured\_simplex\_pathConnected}\\
Source-domain sign and fixed index
 & \texttt{F.continuous\_family}\\
Full conjunction of the two source questions
 & \texttt{F.mo513954\_revision2}\\
General connected domain; positive-pivot classification;
rank identity; $n=2$ counterexample; countable punctures;
block removal; coefficient changes; sphere results
 & Written proofs above; finite rational instances in \texttt{verify.py}\\
\bottomrule
\end{tabularx}
\end{center}
The actual files are linked here:\\
\href{https://github.com/the-omega-institute/trureturing/blob/0f1f822c673c1903dc6b0a209a23dacebca2aa6c/D5/S3/QuadraticForms/StructuredRankOneInertia.lean}{\texttt{StructuredRankOneInertia.lean}}\\
\href{https://github.com/the-omega-institute/trureturing/blob/0f1f822c673c1903dc6b0a209a23dacebca2aa6c/D5/S3/QuadraticForms/StructuredRankOneFamily.lean}{\texttt{StructuredRankOneFamily.lean}}.
The family declaration uses the source's punctured-simplex hypothesis;
the arbitrary-connected-space formulation is supplied by the written proof.

\begin{center}\small
\begin{tabularx}{\linewidth}{@{}>{\raggedright\arraybackslash}p{.53\linewidth}>{\raggedright\arraybackslash}X@{}}
\toprule
Source item & Coverage\\
\midrule
MO 513954 revision 2: both inertia triples; proposed hyperplane proof;
singular null directions
 & Proved: Theorem~\ref{thm:matrix}; prior proof cited~\cite{user1551}.\\
Same revision: nonzero pivot sign; fixed negative coordinate on finite-puncture domain
 & Proved: Theorem~\ref{thm:family}, Corollary~\ref{cor:simplex}.\\
Cross-post comment 11110130: can a negative coordinate hand off?
 & Refuted-and-corrected: its negative-coordinate sets are clopen.\\
Linked MSE 5146381: two-coordinate reduction under the stated hypotheses,
including increasing odd polynomials with decreasing $h(t)/t$
 & Refuted-and-corrected: Theorem~\ref{thm:sphere}.\\
Same question: a specific diffuse limiting configuration
 & Proved for $h(t)=(3t-t^3)/2$: Theorem~\ref{thm:sphere}.\\
Linked discussion: daw's $h=1$ and Exodd's $h=1/t$
 & Prior examples cited; their infima proved above.\\
Same question: numerator bound and stationarity equation
 & Prior proof cited; regularity corrected in Corollary~\ref{cor:stationarity}.\\
Same question: general classification for $h$ and the author's unspecified polynomial
 & Not covered.\\
\bottomrule
\end{tabularx}
\end{center}
The accompanying \texttt{verify.py} checks exact rational instances of the
inertia, kernel-dimension and congruence identities, using symmetric
elimination rather than approximate eigenvalues. The finite test families
and parameter ranges are specified in \texttt{verify.py}, and their summary
counts are recorded in \texttt{verification.json}; these checks illustrate
the identities and are separate from the proofs for all parameters.

\begin{thebibliography}{9}\raggedright
\bibitem{lucid} Lucid (MathOverflow user 596445).
\emph{Inertia of a singular diagonal matrix after a structured rank-one
downdate}, MathOverflow question 513954, revision 2, August 5, 2026.
\href{https://mathoverflow.net/revisions/513954/2}{Version-specific source}.
The source and proposed hyperplane argument are attributed to Lucid;
the post is licensed CC BY-SA 4.0.
\bibitem{user1551} user1551.
Answer to Lucid's \emph{Inertia of a singular diagonal matrix after a
structured rank-one downdate}, Mathematics Stack Exchange question
5146335, answer 5146353, August 7, 2026.
\href{https://math.stackexchange.com/a/5146353}{Original answer},
licensed CC BY-SA 4.0.
\bibitem{otherquestion} Lucid (Mathematics Stack Exchange user 1773327).
\emph{Lower bound for $\langle x,h(x)\rangle/\|h(x)\|_2$ on the unit
sphere: does it reduce to few active coordinates?}, question 5146381,
August 7, 2026.
\href{https://math.stackexchange.com/questions/5146381}{Source question}.
\bibitem{daw} daw. Comment 11110291 on Lucid's Mathematics Stack Exchange
question 5146381, August 8, 2026.
\href{https://math.stackexchange.com/questions/5146381#comment11110291_5146381}{Original comment}.
The choice $h(t)=1$ is credited to this comment (CC BY-SA 4.0).
\bibitem{exodd} Exodd. Answer 5146715 to the same question,
August 11, 2026.
\href{https://math.stackexchange.com/a/5146715}{Original answer},
licensed CC BY-SA 4.0. The example $h(t)=1/t$ is credited to this answer.
\bibitem{inertialean} Trureturing contributors.
\emph{StructuredRankOneInertia.lean}, commit
\texttt{0f1f822c673c1903dc6b0a209a23dacebca2aa6c}.
Pinned source linked in Section~\ref{sec:coverage}.
\bibitem{familylean} Trureturing contributors.
\emph{StructuredRankOneFamily.lean}, same commit.
Pinned source linked in Section~\ref{sec:coverage}.
\end{thebibliography}
\end{document}
